\scanpage{54}
% upper right, boxed pencil note
\margin{\figure{\begin{tikzpicture}[baseline=(copie.base)]
\node[align=left,inner sep=4pt] (copie) {Faire copie\\pour Jouanolou};
\draw[thin] ([xshift=-3pt,yshift=1pt]copie.north west) .. controls ([xshift=4pt,yshift=1pt]copie.north west) and ([xshift=-5pt,yshift=4pt]copie.north east) .. (copie.north east) -- ([yshift=-4pt]copie.north east);
\draw[thin] ([yshift=-1pt]copie.north west) -- (copie.south west) -- ([xshift=2pt,yshift=1pt]copie.south east) -- ([xshift=4pt,yshift=-12pt]copie.north east);
\end{tikzpicture}}}
% letterhead
\begin{center}
{\sffamily HARVARD UNIVERSITY\\
CAMBRIDGE, MASSACHUSETTS 02138}
\end{center}
% letterhead
\noindent {\sffamily DEPARTMENT OF MATHEMATICS\\
2 DIVINITY AVENUE}\\
\hfill Aug 4, 1969\\
% isolated mark at the far right of the date; function unassigned
\margin{\figure{\begin{tikzpicture}[x=1mm,y=1mm]\fill (0,0) .. controls (-0.35,0.05) and (-0.38,0.36) .. (-0.16,0.40) .. controls (0.01,0.45) and (0.11,0.18) .. (0.29,0.31) .. controls (0.40,0.17) and (0.31,-0.12) .. (0.12,-0.06) -- cycle;\end{tikzpicture}}}

\noindent Dear Grothendieck,\\

I have been reading recently some notes from the IHES (one of your seminars, I suppose) entitled ``Exposé \(\overline{\underline{\mathrm{V}}}\). Systèmes
\reversemarginpar\margin{\figure{\begin{tikzpicture}[x=1cm,y=1cm]
\draw[thin] (0,1.10) -- (-0.01,0.13) .. controls (-0.01,0.06) and (0,0.02) .. (0.02,0);
\end{tikzpicture}}}\normalmarginpar%
projectifs J-adiques, par J.P. Jouanolou.'' Could you tell me (for purposes of reference) to which seminar they belong, and whether there is a more recent edition?

Also I have some comments and questions. The most serious question is that I do not understand the proof of the ``Theorem of Shih'', \unsure{A.3.1|A.3.1.} on page 46. It seems to me that the lemma on p. 45 is valid only for \(r>p\) (cf EGA \(\unsure{0|O}_{\overline{\underline{\mathrm{III}}}}\) 13.5.5.2), and so \deleted{\unsure{the|from}} the equivalence of statements (i) and (ii) of the theorem is not clear. (I do agree, however, with the weaker statement ``(resp.\ldots)'' in parentheses). This casts doubt on the following theorem A.3.2 on p. 47, and also on the earlier Prop. 5.3.1 on p. 40. These difficulties would disappear if one worked \unsure{throughout|throughtout} in the category \(\mathrm{pro}\text{-}A\) \margin{\figure{\unsure{\begin{tikzpicture}[x=1mm,y=1mm]
\fill (0,0) .. controls (-0.28,0.02) and (-0.25,0.24) .. (-0.08,0.35) .. controls (0.08,0.43) and (0.26,0.10) .. (0.20,0.03) -- cycle;
\fill (0.55,0.92) .. controls (0.30,0.83) and (0.35,1.13) .. (0.48,1.27) .. controls (0.59,1.48) and (0.81,1.20) .. (0.76,1.06) -- cycle;
\end{tikzpicture}}}} of projective systems with the usual equivalence, rather than in the \deleted{category} finer category \(P_{AR}\).

Other comments: p. 42 line \(\dot{-}1\). Isn't the spectral sequence actually biregular in case of a finite filtration?

p. 47 line \(-1\). I think the hypothesis should be stated for \(n\) and \(n+1\), instead of \(n\) and \(n-1\) (cf EGA \(0_{\overline{\underline{\mathrm{III}}}}\) 11.1.10).
% archive sheet number: 53
\clearpage
